% GENERATED FILE. Edit canonical lessons, then run npm run generate:books.
% canonical-source-hash: fnv1a64:c02b008d9d204e26
% canonical-lessons: JA-C318-setsumei-suru, JA-C318-nattoku-suru, JA-C318-konkyo, JA-C318-imi, JA-C318-naze

\chapter{Explaining Why}
\label{ch:ja-explaining-why}
\hlchaptermodality{318}

\begin{chapteropening}
\textbf{By the end of this chapter:} \emph{I can explain my reasons and ask for an explanation.}

Complete the last lesson of chapter 318: I can explain my reasons and ask for an explanation.
\end{chapteropening}

\section[\texorpdfstring{\ja{せつめいする} (setsumei suru)}{せつめいする (setsumei suru)}]{\ja{せつめいする} (setsumei suru) — to explain}

\label{lesson:JA-C318-setsumei-suru}



\begin{quote}
Before the new one: say the Japanese for to propose, then the Japanese for a conclusion.
\end{quote}

\subsection*{You'll want to know: \ja{せつめいする}}
\textbf{\ja{せつめいする}} — \emph{setsumei suru} — \textquotedblleft{}to explain\textquotedblright{}.

\begin{quote}
\textbf{\ja{わかりませんから、もういちど} \ja{せつめいして} \ja{ください。}} — \emph{wakarimasen kara, mō ichido setsumei shite kudasai} — \textquotedblleft{}I don't understand, so please explain it once more.\textquotedblright{}
\end{quote}

Polite \textbf{\ja{せつめいします}}, \emph{setsumeishimasu}; a \ja{する} verb. Past (polite): \textbf{\ja{せんせいが} \ja{せつめいしました。}}, \emph{sensei ga setsumei shimashita}.

\begin{cousinweb}
\textbf{\ja{せつめいする}} (\emph{setsumei suru}) is built from a Sino-Japanese noun and the native verb \emph{suru}, \textquotedblleft{}to do\textquotedblright{}. The noun, \emph{setsumei}, is two Chinese characters read with their on'yomi, readings borrowed from Chinese: \emph{setsu}, \textquotedblleft{}explanation\textquotedblright{} (Middle Chinese \emph{sywet}) and \emph{mei}, \textquotedblleft{}bright\textquotedblright{} (\emph{mjaeng}) (KANJIDIC2; JMdict; Middle Chinese in Baxter's transcription).
\end{cousinweb}

\begin{grammarlens}[title={What you've built}]
One more word for giving reasons and opinions.
\end{grammarlens}

\subsection*{Your turn}
\begin{itemize}
\raggedright
  \item \emph{Say it:} \emph{setsumei suru}
  \item \emph{Say it:} \emph{setsumei suru}, once more
  \item \emph{Say it:} \emph{wakarimasen kara, mō ichido setsumei shite kudasai} — I don't understand, so please explain it once more.
  \item \emph{Recall:} say the Japanese for to propose, then the Japanese for a conclusion, then say \emph{setsumei suru} again
\end{itemize}

\begin{tcolorbox}[breakable,colback=teal!4,colframe=teal!35!black,title={Before you move on}]
What does \ja{せつめいする} mean? (\textquotedblleft{}To explain\textquotedblright{}.) Say it once more.
\end{tcolorbox}
\section[\texorpdfstring{\ja{なっとくする} (nattoku suru)}{なっとくする (nattoku suru)}]{\ja{なっとくする} (nattoku suru) — to be convinced; to accept a reason}

\label{lesson:JA-C318-nattoku-suru}



\begin{quote}
Before the new one: say the Japanese for a conclusion, then the Japanese for to explain.
\end{quote}

\subsection*{You'll want to know: \ja{なっとくする}}
\textbf{\ja{なっとくする}} — \emph{nattoku suru} — \textquotedblleft{}to be convinced; to accept a reason\textquotedblright{}.

\begin{quote}
\textbf{\ja{りゆうが} \ja{わかったので、なっとくしました。}} — \emph{riyū ga wakatta node, nattoku shimashita} — \textquotedblleft{}Once I understood the reason, I was convinced.\textquotedblright{}
\end{quote}

Polite \textbf{\ja{なっとくします}}, \emph{nattokushimasu}; a \ja{する} verb. Past (polite): \textbf{\ja{ちちは} \ja{なっとくしました。}}, \emph{chichi wa nattoku shimashita}.

\begin{cousinweb}
\textbf{\ja{なっとくする}} (\emph{nattoku suru}) is built from a Sino-Japanese noun and the native verb \emph{suru}, \textquotedblleft{}to do\textquotedblright{}. The noun, \emph{nattoku}, is two Chinese characters read with their on'yomi, readings borrowed from Chinese: \emph{nō}, \textquotedblleft{}settlement\textquotedblright{} (Middle Chinese \emph{nop}) and \emph{toku}, \textquotedblleft{}gain\textquotedblright{} (\emph{tok}) (KANJIDIC2; JMdict; Middle Chinese in Baxter's transcription). The two join as \emph{nat-toku}.
\end{cousinweb}

\begin{grammarlens}[title={What you've built}]
One more word for giving reasons and opinions.
\end{grammarlens}

\subsection*{Your turn}
\begin{itemize}
\raggedright
  \item \emph{Say it:} \emph{nattoku suru}
  \item \emph{Say it:} \emph{nattoku suru}, once more
  \item \emph{Say it:} \emph{riyū ga wakatta node, nattoku shimashita} — Once I understood the reason, I was convinced.
  \item \emph{Recall:} say the Japanese for a conclusion, then the Japanese for to explain, then say \emph{nattoku suru} again
\end{itemize}

\begin{tcolorbox}[breakable,colback=teal!4,colframe=teal!35!black,title={Before you move on}]
What does \ja{なっとくする} mean? (\textquotedblleft{}To be convinced; to accept a reason\textquotedblright{}.) Say it once more.
\end{tcolorbox}
\section[\texorpdfstring{\ja{こんきょ} (konkyo)}{こんきょ (konkyo)}]{\ja{こんきょ} (konkyo) — grounds; a basis (for a claim)}

\label{lesson:JA-C318-konkyo}



\begin{quote}
Before the new one: say the Japanese for to explain, then the Japanese for to be convinced.
\end{quote}

\subsection*{You'll want to know: \ja{こんきょ}}
\textbf{\ja{こんきょ}} — \emph{konkyo} — \textquotedblleft{}grounds; a basis (for a claim)\textquotedblright{}.

\begin{quote}
\textbf{\ja{あの} \ja{いけんには} \ja{こんきょが} \ja{ないと} \ja{おもいます。}} — \emph{ano iken ni wa konkyo ga nai to omoimasu} — \textquotedblleft{}I think that opinion has no grounds.\textquotedblright{}
\end{quote}

\begin{cousinweb}
\textbf{\ja{こんきょ}} (\emph{konkyo}) is Sino-Japanese (kango): two Chinese characters read with their on'yomi, readings borrowed from Chinese. They are \emph{kon}, \textquotedblleft{}root\textquotedblright{} (Middle Chinese \emph{kon}) and \emph{kyo}, \textquotedblleft{}foothold\textquotedblright{} (\emph{kjoH}) (KANJIDIC2; Middle Chinese in Baxter's transcription).
\end{cousinweb}

\begin{grammarlens}[title={What you've built}]
One more word for giving reasons and opinions.
\end{grammarlens}

\subsection*{Your turn}
\begin{itemize}
\raggedright
  \item \emph{Say it:} \emph{konkyo}
  \item \emph{Say it:} \emph{konkyo}, once more
  \item \emph{Say it:} \emph{ano iken ni wa konkyo ga nai to omoimasu} — I think that opinion has no grounds.
  \item \emph{Recall:} say the Japanese for to explain, then the Japanese for to be convinced, then say \emph{konkyo} again
\end{itemize}

\begin{tcolorbox}[breakable,colback=teal!4,colframe=teal!35!black,title={Before you move on}]
What does \ja{こんきょ} mean? (\textquotedblleft{}Grounds; a basis (for a claim)\textquotedblright{}.) Say it once more.
\end{tcolorbox}
\section[\texorpdfstring{\ja{いみ} (imi)}{いみ (imi)}]{\ja{いみ} (imi) — a meaning}

\label{lesson:JA-C318-imi}



\begin{quote}
Before the new one: say the Japanese for to be convinced, then the Japanese for grounds for a claim.
\end{quote}

\subsection*{You'll want to know: \ja{いみ}}
\textbf{\ja{いみ}} — \emph{imi} — \textquotedblleft{}a meaning\textquotedblright{}.

\begin{quote}
\textbf{\ja{いみが} \ja{わからないから、せんせいに} \ja{ききます。}} — \emph{imi ga wakaranai kara, sensei ni kikimasu} — \textquotedblleft{}I don't know what it means, so I will ask the teacher.\textquotedblright{}
\end{quote}

\begin{cousinweb}
\textbf{\ja{いみ}} (\emph{imi}) is Sino-Japanese (kango): two Chinese characters read with their on'yomi, readings borrowed from Chinese. They are \emph{i}, \textquotedblleft{}mind\textquotedblright{} (Middle Chinese \emph{'iH}) and \emph{mi}, \textquotedblleft{}taste\textquotedblright{} (\emph{mj+jH}) (KANJIDIC2; Middle Chinese in Baxter's transcription).
\end{cousinweb}

\begin{grammarlens}[title={What you've built}]
One more word for giving reasons and opinions.
\end{grammarlens}

\subsection*{Your turn}
\begin{itemize}
\raggedright
  \item \emph{Say it:} \emph{imi}
  \item \emph{Say it:} \emph{imi}, once more
  \item \emph{Say it:} \emph{imi ga wakaranai kara, sensei ni kikimasu} — I don't know what it means, so I will ask the teacher.
  \item \emph{Recall:} say the Japanese for to be convinced, then the Japanese for grounds for a claim, then say \emph{imi} again
\end{itemize}

\begin{tcolorbox}[breakable,colback=teal!4,colframe=teal!35!black,title={Before you move on}]
What does \ja{いみ} mean? (\textquotedblleft{}A meaning\textquotedblright{}.) Say it once more.
\end{tcolorbox}
\section[\texorpdfstring{\ja{なぜ} (naze)}{なぜ (naze)}]{\ja{なぜ} (naze) — for what reason; why}

\label{lesson:JA-C318-naze}



\begin{quote}
Before the new one: say the Japanese for grounds for a claim, then the Japanese for a meaning.
\end{quote}

\subsection*{You'll want to know: \ja{なぜ}}
\textbf{\ja{なぜ}} — \emph{naze} — \textquotedblleft{}for what reason; why\textquotedblright{}.

\begin{quote}
\textbf{\ja{なぜ} \ja{はんたいしますか。}} — \emph{naze hantai shimasu ka} — \textquotedblleft{}Why are you against it?\textquotedblright{}
\end{quote}

\begin{cousinweb}
\textbf{\ja{なぜ}} (\emph{naze}) is a native Japanese word (wago), inherited rather than borrowed from Chinese, and its older origin is not settled. JMdict notes that it is usually written in kana (JMdict).
\end{cousinweb}

\begin{grammarlens}[title={What you've built}]
One more word for giving reasons and opinions.
\end{grammarlens}

\subsection*{Your turn}
\begin{itemize}
\raggedright
  \item \emph{Say it:} \emph{naze}
  \item \emph{Say it:} \emph{naze}, once more
  \item \emph{Say it:} \emph{naze hantai shimasu ka} — Why are you against it?
  \item \emph{Recall:} say the Japanese for grounds for a claim, then the Japanese for a meaning, then say \emph{naze} again
\end{itemize}

\begin{tcolorbox}[breakable,colback=teal!4,colframe=teal!35!black,title={Before you move on}]
What does \ja{なぜ} mean? (\textquotedblleft{}For what reason; why\textquotedblright{}.) Say it once more.
\end{tcolorbox}
